\documentclass[12pt]{article}
\usepackage[utf8]{inputenc}
\usepackage{amsmath,amssymb,amsthm}
\usepackage[margin=1in]{geometry}

\newtheorem{theorem}{Theorem}
\newtheorem{lemma}[theorem]{Lemma}
\newtheorem{proposition}[theorem]{Proposition}
\newtheorem{definition}[theorem]{Definition}

\title{Problem 50: Consecutive Prime Sum}
\author{Project Euler Solution}
\date{}

\begin{document}
\maketitle

\section{Problem Statement}

Which prime, below one million, can be written as the sum of the most consecutive primes?

\section{Formal Development}

\begin{definition}[Consecutive Prime Sum]
Let $p_1 = 2 < p_2 = 3 < p_3 = 5 < \cdots$ be the primes in ascending order. A \emph{consecutive prime sum} of length $L$ starting at index~$i$ is
\[
S(i, L) = \sum_{k=i}^{i+L-1} p_k.
\]
\end{definition}

\begin{theorem}[Prefix Sum Formulation]\label{thm:prefix}
Define $P_0 = 0$ and $P_n = \sum_{k=1}^{n} p_k$. Then $S(i, L) = P_{i+L-1} - P_{i-1}$, computable in $O(1)$ after $O(\pi(N))$ preprocessing.
\end{theorem}

\begin{proof}
Telescoping: $S(i,L) = \sum_{k=i}^{i+L-1} p_k = P_{i+L-1} - P_{i-1}$.
\end{proof}

\begin{lemma}[Maximum Feasible Length]\label{lem:maxlen}
Let $L_{\max} = \max\{L : P_L < N\}$. For all $L > L_{\max}$ and all $i \geq 1$, $S(i, L) \geq N$.
\end{lemma}

\begin{proof}
$S(1, L) = P_L$ is the minimum possible sum of $L$ consecutive primes (starting from the smallest). If $P_L \geq N$, then $S(i, L) \geq S(1, L) \geq N$ for all $i$.
\end{proof}

\begin{theorem}[Correctness of Descending Search]\label{thm:correct}
Searching lengths $L = L_{\max}, L_{\max}-1, \ldots, 1$ and returning the first prime $S(i,L) < N$ found yields the prime with the longest consecutive prime sum below $N$.
\end{theorem}

\begin{proof}
Since lengths are tested in decreasing order, the first prime found has maximal $L$. By Lemma~\ref{lem:maxlen}, no length $> L_{\max}$ is feasible. For $L = 1$, every prime $p_i < N$ is a valid sum of length~$1$, so termination is guaranteed.
\end{proof}

\begin{theorem}\label{thm:answer}
The prime below $10^6$ expressible as the longest consecutive prime sum is $997651$, the sum of $543$ consecutive primes starting from $p_4 = 7$.
\end{theorem}

\begin{proof}
Applying the Sieve of Eratosthenes yields $\pi(10^6) = 78\,498$ primes. The prefix sums give $L_{\max} = 546$. The descending search finds that for $L = 543$ and starting index $i = 4$ (i.e., $p_4 = 7$):
\[
S(4, 543) = P_{546} - P_3 = 997651.
\]
Primality of $997651$ is confirmed by the sieve. No lengths $544$, $545$, or $546$ produce a prime sum below $10^6$.
\end{proof}

\section{Complexity Analysis}

\begin{proposition}[Time]
$O(N \log \log N)$ for the sieve, plus $O(\pi(N)^2)$ worst-case for the search (with early termination reducing this substantially in practice).
\end{proposition}

\begin{proof}
The sieve is $O(N \log \log N)$. The nested loop examines at most $\sum_{L} (\pi(N) - L + 1) = O(\pi(N)^2)$ pairs, each costing $O(1)$ via sieve lookup. Early termination upon finding the first prime dramatically reduces the practical cost.
\end{proof}

\begin{proposition}[Space]
$O(N)$ for the sieve; $O(\pi(N))$ for primes and prefix sums.
\end{proposition}

\section{Answer}

\[
\boxed{997651}
\]

\end{document}
